\documentclass[letterpaper,hyper]{HDF5_RFC}

\usepackage[T1]{fontenc}

\renewcommand{\texttt}[1]{\begingroup\small\ttfamily\hyphenchar\font=45
  \spaceskip=.5em plus .3em minus .2em #1\endgroup}
\let\oldtextunderscore\_
\renewcommand{\_}{\oldtextunderscore\allowbreak}

\usepackage{enumitem}
\setlistdepth{4}

\setminted[c]{fontsize=\scriptsize,breaklines,frame=leftline,framesep=2mm,%
  rulecolor=rulecolor,linenos,numbersep=5pt}
\setminted[python]{fontsize=\scriptsize,breaklines,frame=leftline,framesep=2mm,%
  rulecolor=rulecolor,linenos,numbersep=5pt}
\setminted[bash]{fontsize=\scriptsize,breaklines}
\setminted[text]{fontsize=\scriptsize,breaklines}

\title{RFC: String-Based Filter Configuration API\\[4pt]
       {\large Community Review Edition}}
\author{[Scot Breitenfeld]}
\date{2026-05-11}
\rfcversion{2026-001-CR}
\organization{HDFG}
\orglogo{THG_LOGO}

\revision{v1-CR}{May 11, 2026}{Community review edition derived from
  full RFC (tagged \texttt{v1.0-full}).  Technical detail reduced by
  approximately half; focus on user-facing behaviour, compatibility
  guarantees, and open questions for community input.}

\makeindex[intoc]
\addbibresource{support/rfc.bib}

\begin{document}
\maketitle

%% -----------------------------------------------------------------
\begin{abstract}
This RFC proposes a string-based configuration API for HDF5 I/O filters.
The existing \texttt{H5Pset\_filter} interface requires users to construct raw
\texttt{cd\_values} arrays of \texttt{unsigned int}, which forces error-prone
type-punning for floating-point parameters and produces configurations that are
opaque to command-line tools and high-level language bindings.

The proposed design introduces \texttt{H5Pappend\_filter}, which accepts a
filter ID and either a raw \texttt{cd\_values} array or a TOML inline table
parameter string~\cite{toml10}. String parameters are resolved into standard
\texttt{cd\_values} at configuration time; the on-disk format is
\textbf{unchanged}. Filters opt in via an updated plugin class
(\texttt{H5Z\_class3\_t}) with two optional callbacks. Existing call sites
compile and function without modification.

This edition is intended for community review. The full technical specification
is available in the tagged release \texttt{v1.0-full}.
\end{abstract}

\vfill
{\footnotesize\textbf{Acknowledgements:} This work is supported by the U.S.\
Department of Energy, Office of Science, Advanced Scientific Computing Research,
under Contract DE-AC02-06CH11357 and by the U.S.\ Department of Energy, Office
of Science, Office of Advanced Scientific Computing Research, Scientific
Discovery through Advanced Computing (SciDAC) program.}

\makerevisions
\tableofcontents

%% =================================================================
\section{Problem Statement}

Configuring an HDF5 filter today requires constructing a \texttt{cd\_values}
array of \texttt{unsigned int}. This creates three concrete pain points.

\subsection{Type Punning for Floating-Point Parameters}

Filters such as H5Z-ZFP accept \texttt{double} parameters (rate, precision,
tolerance). Callers must bit-cast the value into one or more
\texttt{unsigned~int} slots:

\begin{minted}{c}
unsigned int cd_vals[4];
double rate = 3.5;
cd_vals[0] = H5Z_ZFP_MODE_RATE;
memcpy(&cd_vals[2], &rate, sizeof(double));
H5Pset_filter(plist, H5Z_FILTER_ZFP, flags, 4, cd_vals);
\end{minted}

This is endian-sensitive and essentially impossible to express naturally from
Python, Java, or MATLAB without a dedicated wrapper. In Python the same
operation currently requires:

\begin{minted}{python}
import struct, h5py
rate = 3.5
cd = (1, 0) + struct.unpack("2I", struct.pack("d", rate))
ds = f.create_dataset("x", data=data, compression=32013, compression_opts=cd)
\end{minted}

Every binding (h5py, h5pyd, PyHDF5) reimplements this translation
independently.

\subsection{Opaque CLI Usage}

\texttt{h5repack} users must pass raw integer sequences and rely on external
helper scripts to translate human-readable settings into those sequences.
Example for ZFP rate~3.5:

\begin{minted}{bash}
h5repack -f 'UD=32013,0,4,1,0,0,1074921472' input.h5 output.h5
\end{minted}

\subsection{Opaque Introspection}

The \texttt{cd\_values} array stored in a file is meaningless without
consulting the filter's documentation. There is no standard way to recover
``ZFP rate mode, rate=3.5'' from a stored integer array.

%% =================================================================
\section{Proposed Design}

\subsection{\texttt{H5Pappend\_filter} — New Configuration Entry Point}

\begin{minted}{c}
herr_t H5Pappend_filter(hid_t              plist_id,
                         H5Z_filter_t       id,
                         unsigned           flags,
                         const H5Z_params_t *params);
\end{minted}

\texttt{H5Z\_params\_t} is a tagged union. Two convenience macros cover the
common cases:

\begin{minted}{c}
/* String form */
H5Pappend_filter(plist, H5Z_FILTER_DEFLATE, H5Z_FLAG_MANDATORY,
                 &H5Z_PARAMS_STR("level = 6"));

H5Pappend_filter(plist, H5Z_FILTER_ZFP, H5Z_FLAG_MANDATORY,
                 &H5Z_PARAMS_STR("mode = \"rate\", rate = 3.5"));

/* Blosc2 — sub-compressor specified via dotted keys */
H5Pappend_filter(plist, H5Z_FILTER_BLOSC, H5Z_FLAG_MANDATORY,
                 &H5Z_PARAMS_STR("compressor.name = \"zlib\","
                                 " compressor.level = 6, shuffle = 1"));

/* Parameterless filter */
H5Pappend_filter(plist, H5Z_FILTER_SHUFFLE, H5Z_FLAG_MANDATORY, NULL);

/* Raw cd_values — identical to H5Pset_filter */
unsigned level = 6;
H5Pappend_filter(plist, H5Z_FILTER_DEFLATE, H5Z_FLAG_MANDATORY,
                 &H5Z_PARAMS_RAW(1, &level));
\end{minted}

Equivalent Python (once h5py adopts the new API):

\begin{minted}{python}
ds = f.create_dataset("x", data=data, compression=32013,
                      filter_params="mode = 'rate', rate = 3.5")
\end{minted}

\noindent\textbf{Eager plugin loading.}  Unlike \texttt{H5Pset\_filter}, which
defers all plugin interaction, \texttt{H5Pappend\_filter} with a string
parameter loads the plugin and invokes \texttt{set\_config} at call time.  A
missing plugin returns \texttt{H5E\_NOFILTER} at the call site rather than at
\texttt{H5Dcreate}.

\subsection{\texttt{H5Pget\_filter\_params\_by\_idx} — Human-Readable Introspection}

\begin{minted}{c}
herr_t H5Pget_filter_params_by_idx(hid_t plist_id, unsigned filter_idx,
                                    char *buf, size_t bufsz, size_t *len);
\end{minted}

Returns the human-readable parameter string for the filter at position
\texttt{filter\_idx}.  For filters that implement \texttt{get\_config}, that
output is used; otherwise \texttt{cd\_values} are formatted as a fallback
(e.g., \texttt{"cd\_values=1:0:0:1074921472"}).

%% =================================================================
\section{Parameter String Format}

The parameter string format is a subset of the \textbf{TOML v1.0.0 inline
table syntax}~\cite{toml10}, parsed via a vendored subset of the MIT-licensed
tomlc17 library~\cite{tomlc17}.  Using an established grammar provides typed
values and a stable specification the library does not need to own.

Outer braces are \emph{optional}; the library accepts both forms:

\begin{minted}{text}
level = 6          (bare form — recommended)
{level = 6}        (braced form — equivalent)
\end{minted}

\subsection{Examples}

\begin{center}
\begin{tabularx}{\textwidth}{>{\ttfamily\raggedright\arraybackslash\hsize=1.55\hsize}X
                             >{\raggedright\arraybackslash\hsize=0.45\hsize}X}
\toprule
\textbf{Parameter String} & \textbf{Use Case} \\
\midrule
level = 6 & deflate (integer) \\
mode = "rate", rate = 3.5 & ZFP (string + float) \\
pixels\_per\_block = 16, coding = "entropy" & SZIP \\
scale\_type = "float\_dscale", scale\_factor = 4 & scale-offset \\
fast\_mode = true, level = 6 & boolean flag \\
path = "/data/run\_1,v2/dict.bin" & string value containing a comma \\
compressor.name = "zlib", compressor.level = 6, shuffle = 1
  & Blosc2: sub-compressor (dotted keys) \\
compressor = \{name = "zlib", level = 6\}, shuffle = 1
  & Blosc2: sub-compressor (inline table; equivalent) \\
NULL & no parameters (shuffle, fletcher32) \\
\bottomrule
\end{tabularx}
\end{center}

\subsection{Value Types and Constraints}

Supported value types: integer, float, boolean (\texttt{true}/\texttt{false}),
double-quoted string, single-quoted literal string, nested inline table.
Dotted keys (\texttt{compressor.level}) are equivalent to one level of inline
table nesting and are the recommended form for readability.

Additional constraints imposed by HDF5 on top of TOML:
\begin{itemize}
\item \texttt{inf}, \texttt{-inf}, and \texttt{nan} are valid TOML but
  rejected with \texttt{H5E\_BADVALUE}.
\item Maximum string length: \texttt{H5Z\_CONFIG\_STRING\_MAX}~=~4096~bytes.
\item Maximum parameter count: \texttt{H5Z\_CONFIG\_MAX\_PARAMS}~=~64
  top-level keys.
\item Decimal separator is always \texttt{.} regardless of locale.
\end{itemize}

%% =================================================================
\section{Built-in Filter Parameters}

All built-in filters implement the string API out of the box.

\begin{center}
\renewcommand{\arraystretch}{1.2}
\begin{tabularx}{\textwidth}{lX}
\toprule
\textbf{Filter} & \textbf{Example parameter string} \\
\midrule
Deflate      & \texttt{level = 6} \quad (integer 0--9; default 6) \\
Shuffle      & \texttt{NULL} \\
Fletcher32   & \texttt{NULL} \\
SZIP         & \texttt{coding = "entropy", pixels\_per\_block = 8}
               \quad (\texttt{coding}: \texttt{"entropy"} or \texttt{"nn"};
               ppb: even integer 2--32) \\
N-Bit        & \texttt{NULL} \quad (parameters derived from datatype) \\
Scale-Offset & \texttt{scale\_type = "float\_dscale", scale\_factor = 4}
               \quad (\texttt{scale\_type}: \texttt{"float\_dscale"},
               \texttt{"float\_escale"}, or \texttt{"int"}) \\
\bottomrule
\end{tabularx}
\end{center}

%% =================================================================
\section{For Plugin Authors}
\label{sec:plugin-authors}

Upgrading is \textbf{optional}.  Existing v2 plugins continue to load and
function without any changes.

\subsection{New Plugin Class: \texttt{H5Z\_class3\_t}}

\begin{minted}{c}
typedef struct H5Z_class3_t {
    int                    version;        /* H5Z_CLASS3_T_VERS */
    H5Z_filter_t           id;
    unsigned               encoder_present;
    unsigned               decoder_present;
    const char            *canonical_name; /* required, e.g. "zfp" */
    H5Z_can_apply_func_t   can_apply;
    H5Z_set_local_func_t   set_local;
    H5Z_func_t             filter;
    H5Z_set_config_func_t  set_config;    /* optional: string -> cd_values */
    H5Z_get_config_func_t  get_config;    /* optional: cd_values -> string */
} H5Z_class3_t;
\end{minted}

\subsection{\texttt{set\_config} Callback}

Translates the parameter string into \texttt{cd\_values}. Typed accessor
helpers in \texttt{H5Zdevelop.h} handle parsing; no custom string parsing
is required:

\begin{minted}{c}
static herr_t
my_set_config(const char *params, unsigned *flags,
              size_t *cd_nelmts, unsigned cd_values[], size_t cd_values_size)
{
    int64_t level = 6;   /* default */
    if (params)
        H5Zconfig_get_int(params, "level", &level);

    *cd_nelmts = 1;
    if (cd_values)
        cd_values[0] = (unsigned)level;
    return SUCCEED;
}
\end{minted}

Available accessors: \texttt{H5Zconfig\_get\_int},
\texttt{H5Zconfig\_get\_double}, \texttt{H5Zconfig\_get\_bool},
\texttt{H5Zconfig\_get\_str}.  Dotted keys (e.g.,
\texttt{"compressor.level"}) reach into nested inline tables, supporting
Blosc-style sub-compressor configuration without custom parsing.

\subsection{\texttt{get\_config} Callback}

Reconstructs a parameter string from \texttt{cd\_values} for
\texttt{h5dump~-p} output and \texttt{H5Pget\_filter\_params\_by\_idx}.
Optional but recommended.

\subsection{\texttt{canonical\_name}}

A short, stable string identifier for the plugin class (e.g.,
\texttt{"zfp"}, \texttt{"blosc2"}).  Required for v3 plugins, must not
be \texttt{NULL}, and must not vary between versions of the same plugin.
The library uses \texttt{canonical\_name} as both the programmatic
identifier and the human-readable label returned by
\texttt{H5Pget\_filter2}.  At filter-add time the library writes
\texttt{canonical\_name} into the filter-pipeline message's name field
--- the same on-disk slot v2 plugins have always used to carry their
\texttt{name} --- so the name persists in the file and tools can
display it even when the plugin is not installed.

\noindent\textbf{Compatibility note:} \texttt{cd\_values} is unchanged
by this design; v2 plugins reading files written by a v3-aware library
see exactly the \texttt{cd\_values} their own \texttt{set\_config} (or
the \texttt{H5Pset\_filter} caller) produced.  No \texttt{cd\_nelmts}
surprises for legacy plugins.

%% =================================================================
\section{CLI Tool Changes}

\subsection{\texttt{h5repack}}

The existing syntax continues to work unchanged.  A new form accepts a
filter name and TOML parameter string:

\begin{minted}{bash}
# Old syntax (still works)
h5repack -f 'UD=32013,0,4,1,0,0,1074921472' in.h5 out.h5

# New syntax
h5repack -f 'zfp, mode = "rate", rate = 3.5' in.h5 out.h5
h5repack -f 'deflate, level = 9' in.h5 out.h5
\end{minted}

\subsection{\texttt{h5dump}}

\texttt{h5dump~-p} appends a \texttt{PARAMS\_STRING} line for each filter.
Default output (without \texttt{-p}) is unchanged.

\begin{minted}{text}
FILTERS {
  COMPRESSION DEFLATE { LEVEL 6 }
  PARAMS_STRING "level = 6"
}
\end{minted}

%% =================================================================
\section{Compatibility}

\begin{center}
\renewcommand{\arraystretch}{1.25}
\begin{tabularx}{\textwidth}{lX}
\toprule
\textbf{Area} & \textbf{Status} \\
\midrule
\texttt{H5Pset\_filter} and all existing public APIs
  & Unchanged; no deprecation \\
Existing v1/v2 plugins
  & Load and function without modification \\
On-disk format
  & \textbf{No change} — files written by \texttt{H5Pappend\_filter}
    are byte-identical to equivalent \texttt{H5Pset\_filter} output \\
Files from \texttt{H5Pappend\_filter}
  & Readable by all existing HDF5 readers \\
ABI
  & New functions added; no existing signatures changed
    (minor version bump) \\
h5py, PyTables, netCDF4-Python
  & No changes required; they call \texttt{H5Pset\_filter} \\
Fortran and Java bindings
  & Updated to expose \texttt{H5Pappend\_filter} and
    \texttt{H5Pget\_filter\_params\_by\_idx} \\
\bottomrule
\end{tabularx}
\end{center}

A future RFC targeting HDF5~3.0 will explore typed on-disk parameter storage
(\texttt{H5O\_PLINE\_VERSION\_3}).  The \texttt{cd\_values} packing convention
introduced by this RFC is designed as its foundation; plugins that adopt it now
will be forward-compatible without changes to their callbacks.

%% =================================================================
\newpage
\phantomsection
\addcontentsline{toc}{section}{References}
\printbibliography

\end{document}
